\documentclass[fontset=mac]{ctexart}

% ====== 页面与数学宏包 ======
\usepackage[top=2.54cm, bottom=2.54cm, left=3.18cm, right=3.18cm]{geometry}
\usepackage{amsmath, amssymb, amsthm}
\usepackage{bm}
\usepackage{hyperref}
\usepackage{booktabs}
\usepackage{graphicx}
\usepackage{float}
\usepackage{array}

% ====== 量子算符 ======
\newcommand{\adag}[1]{\hat{a}_{#1}^\dagger}
\newcommand{\aop}[1]{\hat{a}_{#1}}
\newcommand{\Xop}{X}
\newcommand{\Yop}{Y}
\newcommand{\Zop}{Z}
\newcommand{\Iop}{I}
\newcommand{\Nop}[1]{\hat{N}_{#1}}
\newcommand{\ket}[1]{|#1\rangle}
\newcommand{\bra}[1]{\langle #1|}
\newcommand{\braket}[2]{\langle #1|#2\rangle}
\newcommand{\HS}{\mathcal{H}}
\newcommand{\Hred}{\mathcal{H}_{\text{red}}}
\newcommand{\Hphys}{\mathcal{H}_{\text{phys}}}
\newcommand{\Hgauge}{\mathcal{H}_{\text{gauge}}}
\newcommand{\Stabilizer}{\mathcal{S}}

% ====== 标题 ======
\title{Problem 2.1: Build Hamiltonian and Mapping Comparison}
\author{nyloong}
\date{\today}

\begin{document}
\maketitle

% ===========================================================================
\section*{Problem Statement}

将分子哈密顿量放置到量子计算机上，首先需要把费米子哈密顿量映射为 qubit
Pauli 算符。Jordan-Wigner (JW) 和 Parity 是两种常见的映射——它们是互斥的替代方案。
闭壳层分子具有 $N_\alpha$ 和 $N_\beta$ 粒子数守恒对称性（Z$_2$ 对称性），
每种对称性可消除 1 个 qubit。

\textbf{给定分子:} (a) H$_2$ (STO-3G) 和 (d) N$_2$ (STO-3G)。

\textbf{求解:}
\begin{enumerate}
    \item[(a)] 对 H$_2$, 用 PySCF 计算 RHF，提取单电子积分 $h_{pq}$（$p,q$ 为自旋轨道指标）和双电子积分 $(pq|rs)$（化学记号, $p,q,r,s$ 为自旋轨道指标）。写出第二量子化哈密顿量
    \[
    \hat{H} = \sum_{pq} h_{pq} \adag{p} \aop{q}
            + \frac{1}{2} \sum_{pqrs} (pq|rs) \adag{p} \adag{r} \aop{s} \aop{q}
            + E_{\text{core}}.
    \]

    \item[(b)] 以 $h_{01}\adag{0}\aop{1}$（下标 $0,1$ 为自旋轨道指标）为例，推导其在 JW 映射下的 Pauli 串展开式。

    \item[(c)] 统计 H$_2$ 在 JW 和 Parity 两种映射下的 Pauli 项数及最大 Pauli 权重（非 $I$ 的算符数）。

    \item[(d)] 利用 $N_\alpha=1$, $N_\beta=1$ 两个 Z$_2$ 对称性，将 H$_2$ 从 4 个 qubit 约化为 2 个 qubit。显式写出约化哈密顿量；解释 JW 和 Parity 在实现此约化时的区别。

    \item[(e)] 对 N$_2$ (20 个 qubit) 重复 (c) 的统计。

    \item[\textbf{Advanced Challenge}] BK 映射的 Z-串长度为 $O(\log n)$，最大权重远低于 JW。对 N$_2$ ($n=10$ 空间轨道），BK 的最大权重是多少？电路深度低多少倍？
\end{enumerate}

\newpage

% ===========================================================================
\section*{(a) H$_2$/STO-3G 哈密顿量构造}

\subsection*{PySCF 计算流程}

采用键长 $R=0.74$\,\AA, STO-3G 基组:
\begin{itemize}
    \item 空间轨道数: $n_{\text{orb}}=2$ (成键 $\sigma_g$ + 反键 $\sigma_u^*$)
    \item 自旋轨道数: $n_{\text{spin}}=2n_{\text{orb}}=4$ (指标 $0,1,2,3$)
    \item 指标约定: $(0,1)=\alpha$ 自旋, $(2,3)=\beta$ 自旋; $(0,2)$ = 成键, $(1,3)$ = 反键
\end{itemize}

RHF 计算结果:
\begin{align*}
    E_{\text{RHF}} &= -1.1167593074\;\text{Ha}, \\
    \varepsilon_0(\sigma_g) &= -0.57855386\;\text{Ha}, \\
    \varepsilon_1(\sigma_u^*) &= +0.67114349\;\text{Ha}.
\end{align*}

\subsection*{玻恩-奥本海默近似与 $E_{\text{core}}$}

在 Born-Oppenheimer 近似下，原子核被冻结在平衡位置 $R=0.74$\,\AA。核-核 Coulomb 排斥能为常数:
\begin{equation}
    E_{\text{core}} = \frac{Z_1 Z_2}{R} = \frac{1}{0.74 \times 1.889726} \approx 0.715104\;\text{Ha},
\end{equation}
其中 $1.889726$ 是 Bohr 到 \AA 的换算因子。PySCF 直接通过 \texttt{mol.energy\_nuc()} 计算此值。

\subsection*{AO $\to$ MO 积分变换}

原子轨道 (AO) 基下的单电子与双电子积分通过 RHF 轨道系数 $C$ 变换到分子轨道 (MO) 基:
\begin{align}
    h_{pq}^{\text{MO}} &= \sum_{\mu\nu} C_{\mu p} \, h_{\mu\nu}^{\text{AO}} \, C_{\nu q}, \\
    (pq|rs)_{\text{MO}} &= \sum_{\mu\nu\lambda\sigma} C_{\mu p} C_{\nu q} C_{\lambda r} C_{\sigma s} \, (\mu\nu|\lambda\sigma)_{\text{AO}}.
\end{align}

\subsection*{自旋轨道展开}

MO 积分拓展到自旋轨道:
\begin{align}
    h_{2p,\,2q} &= h_{2p+1,\,2q+1} = h_{pq}^{\text{MO}}, \\
    (2p,2q|2r,2s) &= (2p,2q|2r\!+\!1,2s\!+\!1)
                     = (2p\!+\!1,2q\!+\!1|2r,2s) \nonumber\\
                   &= (2p\!+\!1,2q\!+\!1|2r\!+\!1,2s\!+\!1) = (pq|rs)_{\text{MO}} .
\end{align}

由于 H$_2$/STO-3G 使用正则 MO 基 (Fock 矩阵对角化后), 单电子积分矩阵仅有对角元非零:
\begin{equation}
    h_{pq} = \begin{cases}
        -1.2533 & p=q=0,1,2,3, \\
        0       & p\neq q.
    \end{cases}
\end{equation}

\subsection*{第二量子化哈密顿量的显式形式}

将积分代入 $\hat{H} = \sum_{pq} h_{pq} \adag{p} \aop{q}
+ \frac{1}{2} \sum_{pqrs} (pq|rs) \adag{p} \adag{r} \aop{s} \aop{q}
+ E_{\text{core}}$，利用空间与自旋对称性约简后，在程序中得到 \textbf{37 项}费米子算符项。

包含的项类型:
\begin{itemize}
    \item 恒等项: $E_{\text{core}}\, I$ (1 项)
    \item 单电子对角项: $h_{00}\adag{0}\aop{0},\; h_{11}\adag{1}\aop{1},\; h_{22}\adag{2}\aop{2},\; h_{33}\adag{3}\aop{3}$ (4 项)
    \item 双电子项: 来自 8 重对称性简化后共 32 项
\end{itemize}

% ===========================================================================
\section*{(b) JW 映射推导: $h_{01}\adag{0}\aop{1}$ 的 Pauli 展开}

\subsection*{Jordan-Wigner 映射规则}

对于 $n$ 个自旋轨道 (模式), 产生/湮灭算符映射为:
\begin{align}
    \aop{j}   &= \Bigl(\bigotimes_{k=0}^{j-1} \Zop_k\Bigr) \otimes
                 \frac{\Xop_j + i\Yop_j}{2}
                 = Z_0 Z_1 \cdots Z_{j-1} \cdot \sigma_j^+, \label{eq:jw-a} \\
    \adag{j}  &= \Bigl(\bigotimes_{k=0}^{j-1} \Zop_k\Bigr) \otimes
                 \frac{\Xop_j - i\Yop_j}{2}
                 = Z_0 Z_1 \cdots Z_{j-1} \cdot \sigma_j^-. \label{eq:jw-adag}
\end{align}

其中 $\sigma_j^+ = (\Xop_j + i\Yop_j)/2$, $\sigma_j^- = (\Xop_j - i\Yop_j)/2$,
前面的 $Z$ 串用于保持不同模式之间的反对易关系。

\subsection*{具体项展开}

H$_2$/STO-3G 有 4 个自旋轨道 ($j=0,1,2,3$)。对 $j=0$:
\begin{equation}
    \adag{0} = \sigma_0^- = \frac{\Xop_0 - i\Yop_0}{2}.
\end{equation}

对 $j=1$, 式 (\ref{eq:jw-a}) 给出:
\begin{equation}
    \aop{1} = \Zop_0 \otimes \sigma_1^+ = \Zop_0 \cdot \frac{\Xop_1 + i\Yop_1}{2}.
\end{equation}

因此:
\begin{equation}
    \adag{0} \aop{1} = \sigma_0^- \cdot (\Zop_0 \otimes \sigma_1^+)
                      = (\sigma_0^- \Zop_0) \otimes \sigma_1^+.
\end{equation}

\subsection*{计算 $\sigma_0^- \Zop_0$}

利用 Pauli 矩阵乘法规则 $XZ = -iY$, $YZ = iX$:
\begin{align}
    \sigma_0^- \Zop_0 &= \frac{\Xop_0 - i\Yop_0}{2} \cdot \Zop_0
                       = \frac{\Xop_0 \Zop_0 - i\Yop_0 \Zop_0}{2} \nonumber \\
                       &= \frac{(-i\Yop_0) - i(i\Xop_0)}{2}
                       = \frac{-i\Yop_0 + \Xop_0}{2}
                       = \frac{\Xop_0 - i\Yop_0}{2} = \sigma_0^-.
\end{align}

注意 $\sigma_0^-$ 与 $\Zop_0$ 反对易: $\sigma_0^- \Zop_0 = -\Zop_0 \sigma_0^-$,
从而 $\sigma_0^- \Zop_0 = \sigma_0^-$ (吸收号差后等价)。

\subsection*{合并结果}

\begin{equation}
    \adag{0} \aop{1} = \frac{\Xop_0 - i\Yop_0}{2} \cdot \frac{\Xop_1 + i\Yop_1}{2}
                      = \frac{(\Xop_0 - i\Yop_0)(\Xop_1 + i\Yop_1)}{4}.
\end{equation}

展开:
\begin{align}
    \adag{0} \aop{1} &= \frac{1}{4}\bigl(
        \Xop_0\Xop_1 + i\Xop_0\Yop_1 - i\Yop_0\Xop_1 + \Yop_0\Yop_1
    \bigr).
\end{align}

\subsection*{哈密顿量中的实际贡献}

哈密顿量包含 $h_{01}\adag{0}\aop{1} + h_{10}\adag{1}\aop{0}$。
由于单电子积分矩阵是 Hermite 的: $h_{10} = h_{01}^* = h_{01}$ (实数基组)。
而 $\adag{1}\aop{0} = (\adag{0}\aop{1})^\dagger$，
所以:
\begin{align}
    h_{01}\adag{0}\aop{1} + h_{10}\adag{1}\aop{0}
    &= h_{01}\bigl(\adag{0}\aop{1} + \text{h.c.}\bigr) \nonumber \\
    &= h_{01} \cdot \frac{1}{4}\bigl(
        2\Xop_0\Xop_1 + 2\Yop_0\Yop_1
       \bigr) \nonumber \\
    &= \frac{h_{01}}{4}\bigl(\Xop_0\Xop_1 + \Yop_0\Yop_1\bigr).
\end{align}

这给出了两个权重为 2 的 Pauli 项: $\Xop_0\Xop_1$ 和 $\Yop_0\Yop_1$, 系数均为 $h_{01}/4$。

\subsection*{备注: H$_2$ 正则 MO 基的特殊情况}

在正则 MO (canonical MO) 基下, Fock 矩阵是对角的: $h_{01} = h_{10} = 0$。
因此这项对 H$_2$ 的哈密顿量无贡献。
但上述推导是\textbf{一般性的}——对任何包含非对角单电子积分的体系同样适用。
H$_2$ 哈密顿量的非零贡献来自对角项 ($h_{00}$, $h_{11}$) 和双电子积分项。

\subsection*{程序验证 (OpenFermion)}

对 JW 映射的正确性, 我们可以用 $h_{00}\adag{0}\aop{0}$ (非零项) 验证:
\begin{align}
    \adag{0} \aop{0} &= \sigma_0^- \sigma_0^+
                      = \frac{\Xop_0 - i\Yop_0}{2} \cdot \frac{\Xop_0 + i\Yop_0}{2} \nonumber \\
                      &= \frac{1}{4}(\Xop_0^2 + \Yop_0^2 + i[\Xop_0,\Yop_0]) \nonumber \\
                      &= \frac{1}{4}(2I + 2\Zop_0)
                      = \frac{I - \Zop_0}{2}.
\end{align}
因此 $h_{00}\adag{0}\aop{0} = \frac{h_{00}}{2}(I - \Zop_0) = -0.6267\,I + 0.6267\,\Zop_0$,
与 OpenFermion 的 \texttt{jordan\_wigner} 输出一致 $\checkmark$.

% ===========================================================================
\section*{(c) H$_2$ Pauli 项统计: JW / Parity / BK}

\subsection*{项数与最大权重对比}

H$_2$/STO-3G 的费米子哈密顿量 (37 项) 经过三种映射，得到:

\begin{table}[H]
\centering
\begin{tabular}{lccc}
\toprule
\textbf{映射} & \textbf{Pauli 项数} & \textbf{最大权重} & \textbf{权重分布} \\
\midrule
JW     & 15 & 4 & $\{0{:}1,\,1{:}4,\,2{:}6,\,4{:}4\}$ \\
Parity & 15 & 4 & $\{0{:}1,\,1{:}2,\,2{:}6,\,3{:}4,\,4{:}2\}$ \\
BK     & 15 & 4 & $\{0{:}1,\,1{:}3,\,2{:}3,\,3{:}5,\,4{:}3\}$ \\
\bottomrule
\end{tabular}
\caption{H$_2$/STO-3G (4 自旋轨道) 三种映射的 Pauli 统计}
\end{table}

\subsection*{JW 映射下的完整 Pauli 项}

JW 映射将 37 个费米子项压缩为 15 个 Pauli 项（权重 $\le 2$ 的对角项 + $4\times$ 权重 4 的非对角项）:

\begin{equation}
\begin{aligned}
H_{\text{JW}} = &
    -0.09707\, I \\
    &+0.17141\, (Z_0 + Z_1) \;-\; 0.22343\, (Z_2 + Z_3) \\
    &+0.16869\, Z_0Z_1 \;+\; 0.12063\, (Z_0Z_2 + Z_1Z_3) \\
    &+0.16593\, (Z_0Z_3 + Z_1Z_2) \;+\; 0.17441\, Z_2Z_3 \\
    &+0.04530\, (Y_0 X_1 X_2 Y_3 - Y_0 Y_1 X_2 X_3
                   - X_0 X_1 Y_2 Y_3 + X_0 Y_1 Y_2 X_3).
\end{aligned}
\end{equation}

值得注意的是，JW 映射下权重为 4 的项仅 4 个，而权重为 1 的有 4 个、权重为 2 的有 6 个。
对于 H$_2$ 这样的小分子，Parity 和 BK 的最大权重也是 4，因为体系太小还看不到 BK 的 $O(\log n)$ 优势。

% ===========================================================================
\section*{(d) Z$_2$ 对称性约化: 4 qubit $\to$ 2 qubit Tapering}

\subsection*{物理直觉: 粒字数守恒提供 Z$_2$ 对称性}

H$_2$ 基态恰好有 $N_\alpha = 1$ ($\alpha$ 自旋电子数) 和 $N_\beta = 1$ ($\beta$ 自旋电子数)。
这意味着算符 $(-1)^{N_\alpha}$ 和 $(-1)^{N_\beta}$ 与哈密顿量对易,
且在基态所在扇区有确定的本征值 $-1$。

JW 映射下, 每个模式的粒字数算符为 $\hat{n}_j = \adag{j} \aop{j} = (I - \Zop_j)/2$,
因此:
\begin{align}
    \Nop{\alpha} &= \hat{n}_0 + \hat{n}_2
                  = \frac{I - \Zop_0}{2} + \frac{I - \Zop_2}{2}
                  = I - \frac{\Zop_0 + \Zop_2}{2}, \\[4pt]
    \Nop{\beta}  &= \hat{n}_1 + \hat{n}_3
                  = \frac{I - \Zop_1}{2} + \frac{I - \Zop_3}{2}
                  = I - \frac{\Zop_1 + \Zop_3}{2}.
\end{align}
从而:
\begin{align}
    S_\alpha &\equiv (-1)^{\Nop{\alpha}} = \Zop_0 \Zop_2, \\
    S_\beta  &\equiv (-1)^{\Nop{\beta}}  = \Zop_1 \Zop_3.
\end{align}
$S_\alpha$ 和 $S_\beta$ 满足 $S_\alpha^2 = S_\beta^2 = I$, 且 $[S_\alpha, S_\beta] = 0$,
构成阿贝尔 Z$_2 \times$ Z$_2$ 对称群。

在 $N_\alpha = N_\beta = 1$ 扇区, 本征值为:
\begin{equation}
    S_\alpha \ket{\psi} = -\ket{\psi}, \qquad
    S_\beta \ket{\psi} = -\ket{\psi}.
\end{equation}
定义 stabilizer $\bar{S}_\alpha = -S_\alpha = -\Zop_0\Zop_2$,
$\bar{S}_\beta = -S_\beta = -\Zop_1\Zop_3$,
使得在目标扇区内 $\bar{S}_\alpha\ket{\psi} = \bar{S}_\beta\ket{\psi} = +\ket{\psi}$。

\subsection*{稳定子空间分解}

4 量子比特全希尔伯特空间 $\HS = (\mathbb{C}^2)^{\otimes 4}$ 可按 $S_\alpha, S_\beta$ 的本征值分解:
\begin{equation}
    \HS = \bigoplus_{\varepsilon_\alpha, \varepsilon_\beta = \pm 1}
          \HS_{\varepsilon_\alpha, \varepsilon_\beta},
\end{equation}
其中 $\HS_{+1, +1}$ 是 $\bar{S}_\alpha = \bar{S}_\beta = +I$ 的 \textbf{稳定子空间}。
$S_\alpha, S_\beta$ 各有 $\pm 1$ 两个本征值, 故有 $2^2 = 4$ 个子空间,
每个维数 $\dim(\HS) / 4 = 16/4 = 4$。

H$_2$ 的基态和所有两个电子的态都位于 $\HS_{+1,+1}$ (\text{即} $N_\alpha = N_\beta = 1$ 扇区)。
问题变为: 只在 $\HS_{+1,+1}$ 这个 4 维子空间上对角化 $H_{\text{JW}}$,
而不是在 16 维全空间上对角化。

\subsection*{Tapering 算法: 用 Clifford 变换消除 2 个 qubit}

tapering 的核心思想是找到 Clifford 幺正变换 $U$, 将每个 stabilizer 旋转到单个 qubit 的 $Z$ 算符:
\begin{equation}
    U \bar{S}_\alpha U^\dagger = \Zop_{t_1}, \qquad
    U \bar{S}_\beta  U^\dagger = \Zop_{t_2},
\end{equation}
其中 $t_1, t_2$ 是被选中的 \emph{靶 qubit}。

在旋转后的基下:
\begin{itemize}
    \item $\bar{S}_\alpha = +I \;\Longrightarrow\; \Zop_{t_1} = +I$, 即 qubit $t_1$ 固定为 $\ket{0}_{t_1}$
    \item $\bar{S}_\beta  = +I \;\Longrightarrow\; \Zop_{t_2} = +I$, 即 qubit $t_2$ 固定为 $\ket{0}_{t_2}$
\end{itemize}

对所有 Pauli 串 $P$:
\begin{itemize}
    \item $Z_{t_i} \to 1$ (该 qubit 固定为 $\ket{0}$, Z 的本征值为 $+1$)
    \item 对其他包含 $X_{t_i}$ 或 $Y_{t_i}$ 的项: 由于 $\langle 0 | X | 0 \rangle = \langle 0 | Y | 0 \rangle = 0$,
          这些项在子空间上的投影为零, 可消去
    \item qubit $t_i$ 被从有效哈密顿量中移除
\end{itemize}

OpenFermion 的 \texttt{taper\_off\_qubits} 实现了上述过程。输入:
\begin{itemize}
    \item $H_{\text{JW}}$ --- 4 qubit Jordan-Wigner 哈密顿量
    \item $\{\bar{S}_\alpha, \bar{S}_\beta\} = \{-\Zop_0\Zop_2, \;-\Zop_1\Zop_3\}$ --- stabilizer 集合
\end{itemize}

\subsection*{对称性验证}

验证 $S_\alpha, S_\beta$ 与 $H_{\text{JW}}$ 对易:
\begin{align}
    [H_{\text{JW}}, \Zop_0\Zop_2] &= 0, \\
    [H_{\text{JW}}, \Zop_1\Zop_3] &= 0.
\end{align}
这等价于粒字数 $N_\alpha, N_\beta$ 各自守恒——哈密顿量不改变任一自旋方向的电子数。

\subsection*{约化结果}

tapering 后得到 2 qubit 约化哈密顿量 ($\Hred \cong \mathbb{C}^2 \otimes \mathbb{C}^2$, 4 个态):

\begin{equation}
\boxed{
H_{\text{red}} = c_0\, II + c_1\, Z_0 I + c_2\, I Z_1 + c_3\, Z_0 Z_1 + c_4\, Y_0 Y_1,
}
\end{equation}

其中数值系数:
\begin{align}
    c_0 &= -0.3383167378, \\
    c_1 &= -0.3948443634, \\
    c_2 &= -0.3948443634, \\
    c_3 &= +0.0112461572, \\
    c_4 &= -0.1812104620.
\end{align}

(用 \texttt{symmetry\_conserving\_bravyi\_kitaev} 得到的是 $\pm$ 号不同的 XX 形式:
$c_1 = c_2 = +0.3948$, $c_4 = +0.1812$, 相当于 $H \otimes H$ 旋转。)

\subsection*{基态能量守恒}

将 $H_{\text{red}}$ 对角化, 最低本征值:
\begin{equation}
    E_0(H_{\text{red}}) = -1.1372838345\;\text{Ha},
\end{equation}
与 FCI 精确值一致 ($\Delta = 8.88\times 10^{-16}$\,\text{Ha})。
这确认了稳定子空间分解的正确性:
我们消去的恰好是不含任何物理信息的规范自由度, 留下的 4 维空间等价于 2 电子扇区。

\subsection*{JW 与 Parity 映射中对称性处理的对比}

\begin{table}[H]
\centering
\begin{tabular}{l|l|c}
\toprule
\textbf{方法} & \textbf{对称性处理方式} & \textbf{结果 qubit 数} \\
\midrule
JW + tapering & 先映射再手动剪枝 (后处理) & 2 \\
SCBK (Parity-based) & 映射时天然编码粒字数守恒 (内建) & 2 \\
\bottomrule
\end{tabular}
\caption{JW 和 Parity 在对称性约化上的差异}
\end{table}

Parity 映射通过其 binary code 结构, 将粒子数信息编码在 update qubits 中,
使得对称性在映射时就被吸收。JW 则需要额外的 tapering 步骤达到相同效果。
本质上是同一物理的两种代数实现。

% ===========================================================================
\section*{(e) H$_2$ vs N$_2$ 对比: JW / Parity / BK}

使用 OpenFermion-PySCF 自动管线对两种分子进行哈密顿量生成与映射。

\subsection*{H$_2$ / STO-3G ($n_{\text{spin}}=4$)}

\begin{table}[H]
\centering
\begin{tabular}{lccc}
\toprule
\textbf{映射} & \textbf{Pauli 项数} & \textbf{最大权重} \\
\midrule
JW     & 15 & 4 \\
Parity & 15 & 4 \\
BK     & 15 & 4 \\
\bottomrule
\end{tabular}
\caption{H$_2$ --- 体系过小, 三种映射的统计无显著差异}
\end{table}

\subsection*{N$_2$ / STO-3G ($n_{\text{spin}}=20$)}

费米子哈密顿量: \textbf{8269 项}。

\begin{table}[H]
\centering
\begin{tabular}{lccc}
\toprule
\textbf{映射} & \textbf{Pauli 项数} & \textbf{最大权重} \\
\midrule
JW     & 2951 & \textbf{20} \\
Parity & 2951 & 20 \\
BK     & 2951 & \textbf{13} \\
\bottomrule
\end{tabular}
\caption{N$_2$ --- BK 的最大权重显著低于 JW/Parity}
\end{table}

N$_2$ 的 Pauli 权重分布:

\begin{table}[H]
\centering
\small
\begin{tabular}{rccc}
\toprule
\textbf{权重 $w$} & \textbf{JW} & \textbf{Parity} & \textbf{BK} \\
\midrule
 1 & 20 & 2 & 15 \\
 2 & 190 & 52 & 54 \\
 3 & 0 & 36 & 153 \\
 4 & 276 & 408 & 138 \\
 5 & 8 & 16 & 187 \\
 6 & 344 & 289 & 179 \\
 7 & 0 & 42 & \textbf{418} \\
 8 & 400 & 441 & 386 \\
 9 & 4 & 24 & \textbf{437} \\
10 & 588 & 469 & \textbf{465} \\
11 & 0 & 37 & 254 \\
12 & 460 & 442 & 180 \\
13 & 8 & 63 & \textbf{84} \\
14 & 328 & 266 & 0 \\
15 & 0 & 50 & 0 \\
16 & 200 & 171 & 0 \\
17 & 4 & 21 & 0 \\
18 & 108 & 98 & 0 \\
19 & 0 & 15 & 0 \\
20 & 12 & 8 & 0 \\
\midrule
\textbf{合计} & 2951 & 2951 & 2951 \\
\bottomrule
\end{tabular}
\caption{N$_2$/STO-3G Pauli 权重分布对比}
\end{table}

\textbf{关键观察:}
\begin{itemize}
    \item JW 和 Parity 的 Pauli 权重分布覆盖整个 $[1,20]$ 区间，最大权重达到 20；
    \item BK 的权重被限制在 $[1,13]$，完全没有权重 $>13$ 的项；
    \item BK 下大部分项集中在权重 $7$--$11$ 之间（共 1865 项, 占 63\%）；
    \item 这是 BK 的 Fenwick 树结构直接导致的结果：宇称预汇总使得 Z-串长度从 $O(n)$ 降至 $O(\log n)$。
\end{itemize}

% ===========================================================================
\section*{Advanced Challenge: BK 的 $O(\log n)$ 优势}

\textbf{问题:} 对 N$_2$ ($n=10$ 空间轨道, 20 自旋轨道), BK 的最大权重是多少？
电路深度比 JW 低多少倍？

\subsection*{答案}

从 N$_2$/STO-3G 的完整计算:

\begin{table}[H]
\centering
\begin{tabular}{lcc}
\toprule
& \textbf{最大 Pauli 权重} & \textbf{渐进复杂度} \\
\midrule
JW     & 20 & $O(n) = n = 20$ \\
BK     & 13 & $O(\log n)$ \\
\bottomrule
\end{tabular}
\caption{BK 最大权重 = 13}
\end{table}

\begin{equation}
\boxed{\text{电路深度比值} = \frac{\max(\text{JW weight})}{\max(\text{BK weight})}
      = \frac{20}{13} \approx 1.54\times}
\end{equation}

\subsection*{物理解释}

在 Trotter 化的量子电路实现中，每个权重为 $w$ 的 Pauli 项需要 $O(w)$ 个两比特纠缠门
（典型的阶梯型 CNOT 网络需要 $w-1$ 个 CNOT）。\textbf{最大权重项决定了电路深度的主导项}，
因为低权重项可以通过并行化来调度，而高权重项由于涉及大量 qubit，很难被完全并行化。

因此 BK 在 N$_2$ 上实现了约 \textbf{1.5$\times$ 的电路深度缩减}。

\subsection*{为什么是 13 而不是理论下限？}

BK 的单体 Majorana 算符的理论最大权重约为 $2\log_2(n_{\text{spin}}) + 1 \approx 2\cdot 4.32 + 1 \approx 9.6$，
即单个 $\gamma_{2j}$ 或 $\gamma_{2j+1}$ 的权重上界约为 9--11。

然而，双电子积分项 $\adag{p}\adag{r}\aop{s}\aop{q}$ 展开后是 \textbf{4 个 Majorana 算符的乘积}:
\begin{equation}
    \adag{p}\adag{r}\aop{s}\aop{q}
    \propto \gamma_{2p}\gamma_{2r}\gamma_{2s}\gamma_{2q} + \cdots
\end{equation}
当 4 个 Majorana 串的翻转集合 (update set) 重叠最大化时，
权重可以达到 $2 \cdot \max(\text{update}) + 2 \cdot \max(\text{parity}) \approx 13$。

\textbf{BK 的实际最大权重 13 来自多体项的组合效应，而非单体极限。}

\subsection*{随体系增长的展望}

\begin{table}[H]
\centering
\begin{tabular}{ccc}
\toprule
\textbf{体系大小 $n$} & \textbf{JW: $O(n)$} & \textbf{BK: $O(\log n)$}\\
\midrule
H$_2$   (4)   & 4  & 4  \quad (无差异) \\
N$_2$   (20)  & 20 & 13 \quad (1.5$\times$) \\
C$_6$H$_6$ (42) & $\sim$42 & $\sim$16 \quad (2.6$\times$) \\
蛋白质活性位点 (100+) & $\sim$100+ & $\sim$18 \quad (5$\times$+) \\
\bottomrule
\end{tabular}
\caption{BK 的优势随体系增大愈发显著}
\end{table}

BK 的 Fenwick 树结构将 parity-check 信息编码在 $O(\log n)$ 个 qubit 中，
直接缩减了费米子-量子比特变换中的 Z-串长度。在 NISQ 时代电路深度极度受限的情境下，
这种 $O(\log n)$ vs $O(n)$ 的差距对实际可计算体系的大小有决定性影响。

\end{document}
